\section{PETS: uncertainty-aware dynamics for planning}\label{sec:pets}
\textbf{Chua et al., 2018.} Deep Reinforcement Learning in a Handful of Trials
using Probabilistic Dynamics Models~\cite{pets}.
\loc{Sections 4--5; probabilistic ensembles, trajectory sampling, and MPC}

\subsection{Reported mechanism in common notation}
PETS fits bootstrapped probabilistic neural dynamics models and propagates
particles through them to evaluate candidate action sequences. Model differences
represent epistemic uncertainty; each model's predicted distribution represents
transition uncertainty. Control uses sampling-based model predictive control,
executing the first selected action and replanning. This is a model-based
planning method, not policy BPTT through a learned simulator.

For model index $b=1,\ldots,B$, the normalized predictive form is
\begin{align}
\widehat P_{\phi_b}(x'\mid x,a)&=\mathcal N\big(x+\mu_{\phi_b}(x,a),
\Sigma_{\phi_b}(x,a)\big),\\
L_b&=\E_{\mathcal D_b}\left[
(\Delta x-\mu_b)^T\Sigma_b^{-1}(\Delta x-\mu_b)+\log\det\Sigma_b\right].
\end{align}
Constants and a common factor of one half are omitted. $\mathcal D_b$ is a
bootstrap sample; the covariance is diagonal in the neural parameterization.
Particle propagation and the finite planning objective can be written
\begin{align}
\hat x_{k+1}^{(p)}&\sim\widehat P_{\phi_{b_{p,k}}}(\cdot\mid\hat x_k^{(p)},a_k),\\
a_{t:t+H-1}^{*}&\approx\arg\max_{a_{t:t+H-1}}
\frac1P\sum_p\sum_{k=0}^{H-1}\gamma^k r(\hat x_k^{(p)},a_k),\\
a_t^{\rm execute}&=a_t^*.
\end{align}
Resampling ensemble identity across steps and retaining it along a particle
are different trajectory-sampling choices. The latter preserves a coherent
model hypothesis through a rollout.

\subsection{What the formulation exposes}
For an equally weighted ensemble, the law of total variance gives
\begin{equation}\label{eq:ensemblevariance}
\operatorname{Var}(x')=\frac1B\sum_b\Sigma_b
+\frac1B\sum_b(\mu_b-\bar\mu)(\mu_b-\bar\mu)^T.
\end{equation}
The first term describes spread inside a model; the second describes disagreement
between models. A deterministic network with an uncertainty score bolted onto
its output is therefore not equivalent to this construction.

Equation~\eqref{eq:ensemblevariance} is a useful design lens, not a calibration
guarantee. If every model ignores a stair riser, all can agree on the same wrong
prediction. If geometry is missing from the input, variance can instead represent
different unobserved environments producing different outcomes. Measuring that
spread does not tell the controller where the obstacle actually is.

There is also a control distinction. Candidate-action search can evaluate
sequences across different predicted outcomes without using $\partial\Fhat/\partial a$.
A derivative-based policy can fail even when such a planner succeeds, or vice
versa. Comparing them distinguishes model quality from gradient usability.

\subsection{Relevance to adaptation and a discriminating experiment}
Use an ensemble as an uncertainty baseline for the new-geometry problem. On
held-out stair contacts, measure whether disagreement ranks actual prediction
errors and failed control updates. Report reliability curves by contact type,
not just a global correlation. In particular, test examples with unseen riser
height and examples with an unseen collision orientation separately.

An uncertainty-gated controller should be compared with an ungated controller
under the same interaction budget, including any extra samples the gate requests.
The research question is whether uncertainty identifies when additional contact
data is useful. Calling the ensemble ``uncertainty aware'' alone does not answer it.

\section{MBPO: controlling the amount of model-generated experience}\label{sec:mbpo}
\textbf{Janner et al., 2019.} When to Trust Your Model: Model-Based Policy
Optimization~\cite{mbpo}.
\loc{Sections 4--5; practical algorithm and branched-rollout analysis}

\subsection{Reported mechanism in common notation}
MBPO learns an ensemble transition/reward model from real experience, starts
short model rollouts from replay states, and trains a model-free policy using
the generated transitions. The practical implementation uses SAC. Its central
choice is the amount of model usage, rather than differentiating a policy through
many learned transition steps. Real collection and model learning continue online.

Let $\mathcal D_{\rm real}$ and $\mathcal D_{\rm model}$ denote replay sources:
\begin{align}
x_0&\sim\mathcal D_{\rm real},\quad a_k\sim\policy(\cdot\mid\hat x_k),\\
(\hat x_{k+1},\hat r_k)&\sim\widehat P_\phi^{x,r}(\cdot\mid\hat x_k,a_k),
\qquad 0\le k<h_{\rm roll},\\
\mathcal D_{\rm model}&\leftarrow\mathcal D_{\rm model}\cup
\{(\hat x_k,a_k,\hat r_k,\hat x_{k+1})\}.
\end{align}
The superscript $x,r$ explicitly denotes a joint state/reward predictor;
$\widehat P_\phi$ without it continues to mean the state-transition distribution.
A soft actor--critic update on the chosen replay mixture has the standard skeleton
\begin{align}
y_t&=r_t+\gamma c_t\E_{a'\sim\policy}
\left[\min_iQ_{\bar\psi_i}(x_{t+1},a')-\alpha\log\policy(a'\mid x_{t+1})\right],\\
L_Q&=\sum_i\E\big[Q_{\psi_i}(x_t,a_t)-\sg(y_t)\big]^2,\\
L_\pi&=\E_{x,a\sim\policy}\left[\alpha\log\policy(a\mid x)-\min_iQ_{\psi_i}(x,a)\right].
\end{align}
The actor differentiates through the critic's action input, not through a chain
of $\Fhat$ Jacobians. These equations express the estimator distinction; they
do not reproduce every replay or target-update hyperparameter.

\subsection{An original coupling argument for short rollouts}
The following is a review derivation, not MBPO's published policy-improvement
theorem. Fix a policy, environment, and initial-state distribution. For this argument only, use the same
bounded current-state reward $|r(x,a)|\le R_{\max}$ in both systems and assume a
uniform transition error $\TV(P,\widehat P)\le\varepsilon$ on every visited input.
Couple the two trajectories so each step remains matched with probability at
least $1-\varepsilon$. By a union bound,
\begin{equation}
\Pr[x_t\ne\hat x_t]\le\min(1,t\varepsilon).
\end{equation}
When the states disagree, their rewards differ by at most $2R_{\max}$, giving
\begin{equation}\label{eq:coupling}
|J_H-\widehat J_H|\le2R_{\max}\sum_{t=0}^{H-1}\gamma^t\min(1,t\varepsilon).
\end{equation}
If a learned reward also has uniformly bounded error $\varepsilon_r$, add
$\varepsilon_r\sum_t\gamma^t$. The difficult premise is uniform error on visited
inputs; an average replay loss does not provide it. Starting near real states
reduces initial distribution mismatch but does not certify the rest of the rollout.
For deterministic point transitions, total variation can already equal one
whenever the next states differ, however small their Euclidean separation.
Thus this stochastic-transition illustration is not a bound obtained by plugging
a state MSE into Equation~\eqref{eq:coupling}.

\subsection{Relevance to the research question}
After transfer to a stair scene, the trusted horizon may shrink specifically near
edges while remaining longer in free swing. That motivates testing horizon choice
against contact-conditioned validation, rather than automatically retaining the
flat-ground horizon.

An MBPO-style baseline is particularly informative because it can use the same
adapted model without depending on its transition derivatives. If synthetic-data
learning works but BPTT does not, a possible explanation is derivative quality
rather than missing predictive information. The comparison must still match
state/action distributions and count extra simulator interaction fairly.

\section{Dreamer: latent state inference and differentiable imagination}\label{sec:dreamer}
\textbf{Hafner et al., ICLR 2020.} Dream to Control: Learning Behaviors by
Latent Imagination~\cite{dreamer}.
\loc{Sections 2--3, Equations 1--8; Section 4, reconstruction objective}

\subsection{Reported mechanism in common notation}
The original Dreamer learns recurrent latent dynamics from images and rewards,
then learns actions and values inside imagined latent trajectories. Its action
optimization uses differentiable imagined values. This review concerns the
original continuous-control Dreamer, not every later method with that name.
Its original posterior $p_\theta$ becomes $\mathsf q_\phi$ here; original prior
$q_\theta$ becomes $\mathsf p_\phi$; original actor parameters $\phi$ become $\theta$.

An RSSM representation can be expressed as
\begin{align}
m_t&=f_\phi(m_{t-1},\ell_{t-1},a_{t-1}),\\
\ell_t&\sim\mathsf q_\phi(\ell_t\mid m_t,\mathcal E_\phi(o_t)),
\quad\hat\ell_t\sim\mathsf p_\phi(\ell_t\mid m_t),\\
\hat o_t&\sim p_\phi(o_t\mid z_t),\quad
\hat r_t\sim p_\phi(r_t\mid z_t),\quad z_t=(m_t,\ell_t).
\end{align}
In the reconstruction formulation, the minimized negative-ELBO form is
\begin{equation}
L_{\rm WM}=\E\sum_t\left[-\log p_\phi(o_t\mid z_t)-\log p_\phi(r_t\mid z_t)
+\beta\KL\big(\mathsf q_\phi(\ell_t\mid m_t,o_t)\|\mathsf p_\phi(\ell_t\mid m_t)\big)\right].
\end{equation}
The posterior incorporates evidence; the prior must advance without the next
image. Imagined actor training therefore starts at inferred states and then uses
the transition prior. For lambda returns $Y_t$ in Equation~\eqref{eq:lambda},
\begin{equation}
\max_\theta\E\sum_tY_t,\qquad
\min_\psi\E\sum_t\big[V_\psi(z_t)-\sg(Y_t)\big]^2.
\end{equation}
Gradients through $Y_t$ are retained for the actor and stopped for value fitting.
Continuous stochastic variables permit the reparameterization
\begin{equation}
\ell=\mu_\phi+\sigma_\phi\odot\epsilon,\qquad
a=\tanh\big(\mu_\theta(z)+\sigma_\theta(z)\odot\eta\big),
\quad\epsilon,\eta\sim\mathcal N(0,I).
\end{equation}
This displays the sampling-gradient mechanism, not every source-specific
constraint on the networks producing means and scales. The sampled noise is
held fixed in an individual backward evaluation.

\subsection{Why posterior reconstruction is not a rollout test}
The posterior has access to $o_t$. A good reconstruction can arise because the
encoder directly supplies the missing information at every step. During imagined
control, those future observations do not exist. The transition prior, not the
observation-conditioned posterior, must maintain the useful state.

For stair adaptation, evaluate open-loop predictions initialized from one
posterior state, then propagate only actions and the prior. Compare errors at
touchdown and occlusion with posterior reconstruction errors. A gap between the
two identifies a failure of predictive state rather than image decoding alone.

The latent state also need not be a geometrically explicit map. A compact memory
can encode enough terrain information to choose an action while still being
unsuitable for checking foot clearance or contact constraints. Conversely,
forcing all pixels to reconstruct can allocate capacity to appearance details
that are unimportant for control. The useful representation is a testable
task-dependent property, not implied by the name ``world model.''

\subsection{An action-sensitivity test for a latent representation}
One possible review-level diagnostic is to compare two action sequences from
the same inferred state. If they lead to different support outcomes in physics,
the prior should preserve a corresponding distinction in predicted rewards,
geometry-relevant outputs, or controller values. Otherwise the latent model may
represent appearance while erasing the intervention that matters.

This is especially important during transfer. A previously useful latent
compression may discard vertical surfaces because they were irrelevant in the
source scene. Reusing its weights is reasonable as initialization, but freezing
the encoder can restrict the information available to any later controller.

\section{ContactNets: learn contact structure rather than smooth state averages}\label{sec:contactnets}
\textbf{Pfrommer, Halm, and Posa, CoRL 2020.} ContactNets: Learning
Discontinuous Contact Dynamics with Smooth, Implicit Representations~\cite{contactnets}.
\loc{Sections 3--4; Equations 8--17}

\subsection{Reported mechanism in common notation}
ContactNets assumes known contact-free dynamics and learns geometric/contact
quantities from state transitions. Signed distances and contact Jacobians
parameterize discontinuous impulses through a mechanics-based resolver.
Training optimizes a contact-consistency loss over unobserved impulses, rather
than differentiating a next-state MSE through the discontinuous forward solve.
Friction is absorbed into the learned tangential parameterization. Original
contact impulses $\lambda$ are denoted $j$ here, avoiding conflict with lambda returns.

Let $d_{\phi,i}(q)$ be learned signed distance and
\begin{equation}
J^n_{\phi,i}(q)=\nabla_qd_{\phi,i}(q)\,\Gamma(q),\qquad
b^{\rm obs}=M(q)(v'-v)-\Delta t\,f_{\rm free}(x,u).
\end{equation}
The tangential Jacobian $J^\tau_\phi$ is parameterized under the paper's
geometric assumptions; it is not universally the gradient of a scalar distance.
For one observed transition, define the following normalized costs:
\begin{align}
L_1(j)&=\norm{\sum_iJ_{\phi,i}^Tj_i-b^{\rm obs}}^2,\\
L_2(j)&=\sum_i d_{\phi,i}(q')^2\norm{j_i}^2,\\
\tilde v'(j)&=v+M^{-1}\left(\Delta t\,f_{\rm free}+\sum_iJ_{\phi,i}^Tj_i\right),\\
\tilde d'_i(j)&=d_{\phi,i}(q)+\Delta t\,J^n_{\phi,i}(q)\tilde v'(j),\\
L_3(j)&=\sum_i\min(0,\tilde d'_i(j))^2,\\
L_4(j)&=\sum_i\norm{\norm{J^\tau_{\phi,i}v'}j_i^\tau+j_i^nJ^\tau_{\phi,i}v'}^2,\\
L_{\rm CN}(\phi)&=\E_{\mathcal D}\min_{\substack{j_i^n\ge0\\\norm{j_i^\tau}\le j_i^n}}
\sum_{k=1}^{4}L_k(j).
\end{align}
The unit friction cone follows the absorbed-friction convention. Predictive
dynamics are subsequently generated by the contact resolver, not by identifying
the training optimizer's impulses with a universal feedforward impulse network.

\subsection{Interpreting the structure}
$L_1$ explains the observed net contact impulse. $L_2$ discourages impulses
away from surfaces. $L_3$ penalizes predicted penetration, and $L_4$ expresses
a frictional dissipation residual. The observed $q'$ fixes geometric quantities
while the inner optimization searches for compatible impulses. At fixed learned
geometry this inner identification problem has useful convex structure; the
neural parameter-learning problem is not thereby globally convex.

The important representation principle is that smooth geometric functions can
parameterize a discontinuous transition. This differs from requiring a smooth
MLP to directly approximate every state jump. The transition may still have
mode boundaries, and differentiating a resulting controller is a separate issue
from learning a physically plausible contact description.

\subsection{What transfers to our problem, and what does not}
For a G1 with known robot geometry, it may be unnecessary to learn every signed
distance from scratch. A first comparison can supply exact static terrain
geometry and learn only uncertain contact parameters or residual effects. This
isolates the value of mechanics structure from visual reconstruction and
geometry identification.

Known mass and contact-free dynamics are substantive assumptions. If they are
wrong, a learned contact term can compensate for inertia or actuator errors and
lose its intended physical interpretation. A contact-structured model must
therefore be tested under controlled identification errors, not simply declared
more transferable than a state MLP.

The relevant experiment compares additional target data needed by an unstructured
conditioned model and a structured contact model, holding environmental inputs
and known physics constant. A reduction in penetration alone is insufficient:
the model must also support actual successful control updates. ContactNets is
the structural precedent, not evidence that this humanoid adaptation result
has already been demonstrated.
